\documentclass{article}
\PassOptionsToPackage{numbers}{natbib}
\usepackage[final,sglblindworkshop]{neurips_2026}
\usepackage[utf8]{inputenc}
\usepackage[T1]{fontenc}
\usepackage{
  hyperref,url,booktabs,tabularx,amsfonts,amsmath,
  microtype,xcolor,graphicx
}
\hypersetup{hidelinks}
\workshoptitle{Galileo Hackathon 2026}

\makeatletter
\renewcommand{\@noticestring}{
  Galileo Hackathon 2026. Participant research report.
}
\makeatother

\newcommand{\pending}[1]{%
  \textcolor{blue!60!black}{\textbf{Unresolved:} #1}
}

\title{Reflection averaging and period copying for wake forecasts}
\author{
  Team Optimus (Raghav Vijaykannan, Pranav Vijay Krishna, Hrithik Yash)\\
  Team ID: G2-12 \quad Task: pitching-wing-wake
}

\begin{document}
\maketitle

\begin{abstract}
With model weights fixed, reflection averaging lowers validation NMSE
from 0.2094 to 0.1856, improving forecasts in all 12 cases; three times
the training exposure lowers it further to 0.1440.
On the hidden tiers, the averaged model's main failure is frequency
extrapolation. A parameter-free forecast that copies the history forward
by whole forcing periods, gated by how well the history repeats,
lowered the portal score from 0.2876 to 0.1789, and a proxy-tuned soft
blend lowered it to 0.1680. Training with a per-window, variance-normalised
loss (validation 0.1440 to 0.1261) lowered it to 0.1593.
The results support these inference changes on the evaluated cases;
they do not establish the benefit of training augmentation or
superiority over generic averaging.
\end{abstract}

\section{Scientific goal \& initial hypothesis}

We forecast 20 vorticity frames from 10 past frames and the pitching
conditions. We test whether teaching or enforcing reflection symmetry
improves these forecasts. Reflection consistency means that reflecting
the input should reflect the predicted wake; an individual wake need
not be symmetric.

We hypothesized that augmentation would reduce reflection inconsistency
and that averaging would remove it and could improve accuracy.
Lower inconsistency without lower forecast error would contradict
the accuracy hypothesis.

For $y\mapsto-y$, $G$ negates the history's vorticity, reverses its
vertical rows, and negates both pitch-phase sine and cosine.
Frequency $k$ and positive amplitude $\theta_0$ remain unchanged.
$S$ negates and vertically reverses the forecast.
For B's predictor $F_B$, C is
\begin{equation}
F_C(x)=\tfrac12\left[F_B(x)+S F_B(Gx)\right].
\end{equation}
This enforces reflection consistency of the forecast map using two
forecast passes.

The task permits this transformation. The domain is symmetric about
$y=0$; exact grid placement and boundary implementation remain
unaudited.

\begin{figure}[!ht]
\centering
\includegraphics[width=\linewidth]{figures/fig1_setting.pdf}
\caption{
Left: released training and validation conditions.
Right: averaging moves B and B+ to zero measured reflection defect
and lower forecast error. Arrows connect the same weights.
Defect measures disagreement between reflecting a prediction and
predicting a reflected input; it does not measure accuracy.
}
\label{fig:setting}
\end{figure}

\newpage
\section{Research trajectory}

Averaging the organizer checkpoint (A$_{\rm ref}$) changed NMSE only
from 0.2571 to 0.2562, with an interval crossing zero.
We next trained B with reflection augmentation; averaging its
predictions (C) improved all 12 cases.
The matched local run without augmentation, A, did not finish.
We then increased training exposure for B+ and applied averaging
to obtain C+. To check whether gains might reflect averaging more
generally, E averaged A$_{\rm ref}$ and B forecasts.
Fresh train-only reproduction verified B/C.

The portal score of C (0.3620) was far above its validation NMSE.
The hidden set adds amplitude ($\theta_0$ 38--45$^\circ$, T3a) and
frequency ($k$ 7.5--9, T3f) extrapolation tiers beyond the training
range ($k\le7$, $\theta_0\le35^\circ$). We therefore built proxies from
training data: a model trained on $k\le5.5$ and scored on held-out
cases with $k$ 6--7, and one trained on $\theta_0\le30^\circ$ and scored
on 31--35$^\circ$. Each test motivated the next submission
(Table~\ref{tab:portal}).

\section{Final model \& comparisons}

B uses the official width-64 U-Net (1,890,708 parameters), seed 0,
and 40 epochs. Each training example is reflected with probability
$0.5$ using a separate random generator.
B uses 22,680 windows from 810 training cases; B+ uses 68,040,
giving three times the training steps.
Validation is excluded from fitting.
The packaged command records the full optimizer recipe.
C/C+ average forecasts and retain statistical outputs.
E equally averages A$_{\rm ref}$ and B.

All validation scores use the same 48 windows from 12 public-validation
cases. Lower NMSE is better. The official rule clips each lead's NMSE at
4 before averaging. Intervals resample paired case means 2,000 times
(seed 0). They describe case-sampling variability, not training-seed
uncertainty or selection-adjusted performance.

\begin{center}
\begin{tabular}{@{}lrrrrr@{}}
\toprule
Model & NMSE & Final lead & Defect & Fit time (s) & Passes\\
\midrule
A$_{\rm ref}$ & 0.2571 & 0.4105 & 0.1638 & Unknown & 1\\
A$_{\rm avg}$ & 0.2562 & 0.3696 & 0 & Reuses A$_{\rm ref}$ & 2\\
B & 0.2094 & 0.3144 & 0.0505 & 1,068 & 1\\
C & 0.1856 & 0.2754 & 0 & Reuses B & 2\\
B+ & 0.1670 & 0.2373 & 0.0493 & 3,881.6 & 1\\
C+ & 0.1440 & 0.2121 & 0 & Reuses B+ & 2\\
E & 0.1897 & 0.2855 & Unmeasured & A$_{\rm ref}$ + B & 2\\
\bottomrule
\end{tabular}
\end{center}

\noindent\textbf{More training.}
C+ lowers validation NMSE by 0.04160 versus C, with interval
$[-0.04747,-0.03616]$ and all 12 cases improving. It uses three times
the training exposure, so this is an accuracy gain at extra training
cost, not an equal-budget symmetry result.

\begin{figure}[!ht]
\centering
\includegraphics[width=\linewidth]{figures/fig2_accuracy.pdf}
\caption{
C improves B across the forecast horizon and in all 12 cases.
Its advantage over E is unresolved: the interval crosses zero.
Dots show cases, diamonds show mean differences, and bars show
95\% case-bootstrap intervals. Faint curves show individual C cases.
Negative differences favor the first model.
}
\label{fig:accuracy}
\end{figure}

\noindent\textbf{Extrapolation and period copying.}
On the frequency proxy, the averaged model scored 0.99. Clamping the
$k$ input to the training range made it worse (1.04). Shrinking the
forecast toward the mean of the 10 history frames, ramped over the
leads and applied only for out-of-range $k$, gave 0.78; the same
shrinkage hurt on the amplitude proxy and on validation, so it is
gated on $k$.
Hiding $k$ during training did not help extrapolation (1.04--1.05).
This suggests the model has not learned the faster dynamics; the cause
is not the $k$ input channel.
The forcing period $T=2\pi/k$ is known. For $k\ge6.98$, one period fits
inside the 0.9\,$c/U$ history, so the forecast can copy the history
forward by whole periods (linear interpolation in time). A gate measures
how well the history repeats with period $T$ (residual $r$, normalised
by the history's variance). The copy has no trained parameters, so we
scored it on 2,100 training windows from 523 cases with $k$ 3.5--7,
lengthening the history to cover one period. Windows with $r<0.3$ had
mean NMSE 0.059. Periodic windows had error 0.044--0.058 in every $k$
bin, so accuracy does not depend on $k$.
On the frequency proxy, with a CNN that never saw $k>5.5$, the copy beat
the CNN even for aperiodic wakes (0.61 versus 0.99).
For out-of-range $k$, a soft blend (full copy below $r=0.5$, none above
1.2) scored 0.1905 against 0.270 for a hard gate at $r<0.3$.
In range, the hard gate is kept, and validation outputs are unchanged.
Copying commutes with $S$, so reflection consistency is preserved.
A half-period copy with reflection did not help on the amplitude proxy
or on validation. Analog forecasting, using the continuation of the
closest training history, scored 1.02 on validation.

\begin{table}[!ht]
\centering
\begin{tabular}{@{}lllr@{}}
\toprule
Submission & Model & Inference change & Portal score\\
\midrule
C\_from\_B & C & reflection averaging & 0.3620\\
C\_from\_Bplus\_ood & C+ & + out-of-range shrinkage & 0.2876\\
\#95 C\_from\_Bplus\_hybrid & C+ & + period copy, hard gate & 0.1789\\
C\_from\_Bplus\_hybrid2 & C+ & + soft blend for $k>7$ & 0.1680\\
C\_nmse\_hybrid & C$_{\rm nmse}$ & same, metric-weighted training loss & 0.1593\\
\bottomrule
\end{tabular}
\caption{Hidden-tier portal scores (lower is better). Each row changes one
step: the network (C$\to$C+$\to$C$_{\rm nmse}$) or the inference procedure.
The portal's aggregation over tiers is not published.}
\label{tab:portal}
\end{table}

\noindent\textbf{Cost and checks.}
B+ took 3.6 times B's recorded fit time for three times the steps;
timing was not controlled. C needs two forecast passes without
additional training; the period copy adds no fitted parameters.
Reproduced C passed shape, finite-output, NMSE, non-frozen-forecast
and latency checks. Its H100 median batch-1 latency was 1.959 ms
over 50 measurements, with other jobs paused and their GPU contexts
resident. Every submitted package passed the same checks on validation.
The worst-case CPU latency with the copy path was 19 ms.
C+ and its derivatives lack a fresh full training reproduction.

\medskip
\noindent\textbf{Artifact status.}
\pending{
Confirm the final selected submission, its portal ID, ZIP hash and code
revision. A metric-weighted-loss variant (C\_nmse\_hybrid) is pending.
}

\begin{figure}[!ht]
\centering
\includegraphics[width=\linewidth]{figures/fig3_fields_robust.pdf}
\caption{
Early and late forecasts, followed by late absolute error.
Compare vortex locations across B, C and C+.
This example does not replace the case-level results.
Case \texttt{h1\_008}, start 60, was selected nearest the median
reference error, not by improvement.
Shared limits are $\pm26.044$ for vorticity and $[0,13.546]$
for error (99th-percentile limits; extremes saturate).
}
\label{fig:fields}
\end{figure}

\section{Key insights}

\noindent\textbf{Averaging helps this trained model.}
With B's weights fixed, C lowers mean NMSE by 0.02381 (11.4\%).
The 95\% case interval is $[-0.02734,-0.02038]$, with all 12 cases
improving (Fig.~\ref{fig:accuracy}).
This comparison isolates the inference operation on these cases.

\medskip
\noindent\textbf{Consistency does not guarantee accuracy.}
A$_{\rm avg}$ also has zero defect, but its paired change is
$-0.00097$, with interval $[-0.07174,0.05850]$.
C versus E is likewise inconclusive:
$-0.00409$, with interval $[-0.02451,0.01144]$;
5 of 12 cases favor C.
These results establish neither equivalence nor a symmetry-specific
advantage over generic averaging.

\medskip
\noindent\textbf{Known periodicity beats learned extrapolation.}
For frequencies above the training range, copying the history by the
known forcing period outperformed the network on held-out proxy cases.
This held for periodic wakes (0.055 versus 0.585) and for aperiodic
ones (0.61 versus 0.99). Adding the copy lowered the portal score from
0.2876 to 0.1789, and the soft blend lowered it to 0.1680.
The comparison uses proxies built from training data, so the size of
the gain on the hidden tier is inferred from the portal, not measured
directly.

\newpage
\section{Limitations \& next experiment}

Without a completed matched local A, B versus the organizer checkpoint
cannot isolate augmentation.
One seed and repeated use of 12 validation cases limit accuracy claims.
The exact discrete physical reflection remains unaudited.
Portal scores are aggregate hidden-tier results with an unpublished
aggregation. Our extrapolation evidence comes from proxies built on
training data, with shorter training (8--12 epochs) and a smaller
extrapolation distance than the hidden tiers. Thresholds were chosen
on these proxies. The copy applies only when one period fits in the
history ($k\ge6.98$). Amplitude extrapolation remains unaddressed:
clamping $\theta_0$, shrinkage and half-period copying did not help on
the amplitude proxy, and ensembling four models did not help on the
frequency proxy.

Next, complete A with B's initialization, window order, optimizer,
schedule and exposure, changing only augmentation.
Evaluate A, averaged A, B and C on fixed cases, then repeat across seeds.
A matched independently trained predictor is needed to compare
reflection averaging with generic averaging; E is not that control.
For amplitude extrapolation, test training-side changes on the
$\theta_0$ proxy before full runs.

\medskip
\noindent\textbf{Prior work.}
Wang et al.~\citep{wang2021} compare equivariant fluid models with
ordinary and augmented models; forcing and boundaries must respect
the transformation.
Elesedy and Zaidi~\citep{elesedy2021} give squared-risk benefits under
an invariant input distribution and equivariant target, not a
guarantee for our finite, clipped score.
Fromme et al.~\citep{fromme2025} study steerable convection forecasting,
rather than this inference operation.
Our contribution is the measured same-weights wake comparison and the
proxy-validated period copy, not a new method.

\section{Reproducibility, contributions \& AI use}

\noindent\textbf{Evidence.}
Fresh seed-0, 40-epoch train-only reproduction matched B's final
tensors and all C output arrays exactly.
Environment: Python 3.12.16, PyTorch 2.10/CUDA 12.8, H100 80 GB.
The packaged command initializes seed 0, then runs
\texttt{train --mirror --train-only} with only \texttt{train.npz}
visible. Validation is exposed after fitting.
Proxy experiments, period-copy tests and packages are under
\texttt{runs/p3/} and \texttt{submission/} in the repository.

Data:
\href{https://doi.org/10.5281/zenodo.24386760}
{doi:10.5281/zenodo.24386760}.
The
\href{https://github.com/pranavvkrishna/GalileoHackathon-Optimus}
{repository}
and
\href{https://github.com/pranavvkrishna/GalileoHackathon-Optimus/issues/2}
{coordination issue}
identify source and evidence locations.
\pending{Update the reviewed repository snapshot (was \texttt{63d09df}) to the final submission revision.}

\medskip
\noindent Verified artifact:
\texttt{C\_from\_B\_fresh\_verified\_v2.zip}.

\noindent ZIP SHA-256 (concatenate lines):\\
\texttt{0382f3bd83c37a5400b0cc21188df41e30ff}\\
\texttt{6e885b1432aaca43bea52f32bd38}

\noindent C\_from\_Bplus\_hybrid2 ZIP SHA-256:\\
\texttt{9ff281434be68937298866e124c97959}\\
\texttt{5ce1eaccd7ce0d00c21b5a92718877fc}

\medskip
Included plot data and provenance rebuild figures without a GPU.
The C+ example was inferred on CPU; lead errors match the GPU
diagnostic within $2.3\times10^{-7}$.
Full-range plots document color saturation.

\medskip
\noindent\textbf{Human judgment and AI use.}
The team reviewed the results and exercised human judgment in
assessing the scientific claims, limitations and conclusions.
AI assistance supported literature review, code, figures, LaTeX
and drafting. Codex assisted with report preparation;
Undermind supported paper retrieval and reading.
Claude Code assisted with evaluation code, gates, proxy experiments,
the period-copy analysis and submission packaging.
The team reviewed the interpretation and final wording.

\clearpage
\begin{thebibliography}{9}

\bibitem{wang2021}
Rui Wang, Robin Walters, and Rose Yu.
Incorporating Symmetry into Deep Dynamics Models for Improved
Generalization.
\emph{ICLR}, 2021.
\url{https://doi.org/10.48550/arXiv.2002.03061}.

\bibitem{elesedy2021}
Bryn Elesedy and Sheheryar Zaidi.
Provably Strict Generalisation Benefit for Equivariant Models.
\emph{Proceedings of ICML}, PMLR 139, 2021.
\url{https://doi.org/10.48550/arXiv.2102.10333}.

\bibitem{fromme2025}
Fynn Fromme, Christine Allen-Blanchette, Hans Harder,
and Sebastian Peitz.
Surrogate Modeling of 3D Rayleigh--B\'enard Convection with
Equivariant Autoencoders.
ArXiv preprint, 2025.
\url{https://doi.org/10.48550/arXiv.2505.13569}.

\end{thebibliography}
\end{document}
